\chapter{Cơ chế cộng tác của hệ đa tác tử LLM (Tran và cộng sự, 2025)}\label{chap:3}

Bài khảo sát \emph{Multi-Agent Collaboration Mechanisms: A Survey of LLMs} của Tran, Dao, Nguyen, Pham,
O'Sullivan và Nguyen (arXiv:2501.06322, 2025)~\cite{ch3:tran} chuyển trọng tâm từ năng lực của một mô hình ngôn
ngữ lớn (LLM) sang câu hỏi: \emph{nhiều tác tử phối hợp với nhau như thế nào}. Dù bài thí nghiệm không dùng LLM,
khung phân tích của bài khảo sát đủ tổng quát để mô tả cả thị trường hai cửa hàng lẫn hệ mua sắm của Chương~\ref{chap:8}.

\section{Khung phân tích năm chiều}

Bài khảo sát mô tả một hệ cộng tác bằng năm chiều (Hình~\ref{fig:ch3-frame}): \textbf{tác nhân} (actors) --- ai
tham gia; \textbf{kiểu quan hệ} (types) --- hợp tác, cạnh tranh hay vừa hợp tác vừa cạnh tranh;
\textbf{chiến lược} (strategies) --- dựa trên luật, vai trò hay mô hình; \textbf{cấu trúc truyền thông}
(structures) --- tập trung, phi tập trung hay phân cấp; và \textbf{điều phối} (coordination/orchestration) --- tĩnh
hay động. Mỗi chiều vừa là nguồn lợi ích vừa là nguồn chi phí: thêm một tác tử kiểm chứng giảm sai sót nhưng làm
tăng số lời gọi và độ trễ.

\hinhtikz{ch3_frame}{Khung năm chiều mô tả cơ chế cộng tác đa tác tử theo Tran và cộng sự.}{ch3-frame}

\nhanxet{Năm chiều độc lập với nhau: cùng một tập tác tử có thể ghép với nhiều cấu trúc và cách điều phối khác
nhau, nên thiết kế MAS là chọn một điểm trong không gian năm chiều này chứ không chọn một ``kiến trúc mẫu''.}

\section{Kiểu quan hệ: hợp tác, cạnh tranh và cạnh--hợp}

\textbf{Hợp tác} (cooperation) là cùng đóng góp vào một mục tiêu chung, ví dụ nhóm tác tử lập kế hoạch, viết mã và
kiểm thử. \textbf{Cạnh tranh} (competition) là ưu tiên mục tiêu riêng có thể xung đột, ví dụ tranh luận đối kháng để
lộ điểm yếu của lập luận. \textbf{Cạnh--hợp} (coopetition) kết hợp hai dạng trên, như thương lượng tìm thỏa thuận
dưới lợi ích khác nhau~\cite{ch3:tran}. Hình~\ref{fig:ch3-types} đặt hai hệ của báo cáo lên trục này.

\hinhtikz{ch3_types}{Phổ kiểu quan hệ giữa các tác tử và vị trí của hai hệ trong báo cáo.}{ch3-types}

\nhanxet{Thị trường hai cửa hàng là cạnh tranh nhưng \emph{không} tổng bằng không: tổng lợi nhuận thay đổi theo số
đơn và biên lãi, nên cả hai cùng giữ giá có thể tốt hơn cho cả hai so với cùng giảm giá --- đúng dạng thế lưỡng nan
mà Chương~\ref{chap:4} sẽ tính bằng số.}

\section{Chiến lược và cấu trúc truyền thông}

Ở chiều chiến lược, bài khảo sát phân biệt ba loại giao thức. \textbf{Dựa trên luật} (rule-based): hành vi định
sẵn bằng quy tắc, dễ dự đoán và kiểm tra nhưng kém thích nghi. \textbf{Dựa trên vai trò} (role-based): mỗi tác tử
giữ một vai chuyên biệt, hỗ trợ chuyên môn hóa và phân công. \textbf{Dựa trên mô hình} (model-based): tác tử suy luận
xác suất về môi trường và ý định của tác tử khác, xử lý được bất định nhưng đắt và khó triển khai.

\begin{luuy}{Hai nghĩa của ``model-based''}
``Model-based'' trong phân loại cộng tác chỉ việc suy luận về tác tử khác; nó không đồng nghĩa với model-based RL
(học mô hình chuyển trạng thái). Q-learning của bài thí nghiệm là model-free, còn tác tử theo luật là rule-based.
\end{luuy}

Ở chiều cấu trúc, các tác tử có thể giao tiếp qua một nút trung tâm (\textbf{tập trung}), trực tiếp với nhau
(\textbf{phi tập trung}, ngang hàng) hoặc theo nhiều tầng quản lý (\textbf{phân cấp}) (Hình~\ref{fig:ch3-struct}).
Cấu trúc tập trung dễ kiểm soát nhưng có điểm nghẽn; phi tập trung bền với lỗi một nút nhưng tốn thông điệp và khó
đạt đồng thuận; phân cấp cân bằng hai bên cho bài toán lớn. Về điều phối, luồng \textbf{tĩnh} định trước thứ tự các
bước, còn luồng \textbf{động} chọn tác tử và phụ thuộc khi đã biết yêu cầu cụ thể.

\hinhtikz{ch3_struct}{Ba cấu trúc truyền thông: tập trung, phi tập trung và phân cấp.}{ch3-struct}

\section{Ứng dụng, thách thức và ánh xạ vào hai hệ của báo cáo}

Bài khảo sát tổng hợp ứng dụng trong mạng 5G/6G và Công nghiệp 5.0, hỏi đáp và sinh ngôn ngữ, cùng các mô phỏng
xã hội và văn hóa. Kết luận chung là hiệu quả phụ thuộc vào kênh phối hợp, tri thức miền và chi phí mở rộng; trong
một số thiết lập, một tác tử đơn còn tốt hơn nhiều tác tử phối hợp kém. Các vấn đề mở gồm quản trị thống nhất,
ra quyết định tập thể, \textbf{lỗi lan truyền} giữa các tác tử, chi phí tài nguyên và thiếu benchmark chung.
Bảng~\ref{tab:ch3-map} dùng năm chiều để mô tả hai hệ được thực hiện.

\begin{table}[!htbp]\centering\caption{Ánh xạ khung năm chiều vào hai hệ đa tác tử của báo cáo.}\label{tab:ch3-map}
{\small\begin{tabular}{@{}L{2.6cm}L{5.9cm}L{6.0cm}@{}}\toprule
Chiều & Thị trường hai cửa hàng (Ch.~4--7) & Hệ mua sắm (Ch.~8) \\\midrule
Tác nhân & Một tác tử Q-learning, một tác tử theo luật & Catalog, Recommendation, Inventory, Coordinator \\
Kiểu quan hệ & Cạnh tranh lợi nhuận, không tổng bằng không & Hợp tác hoàn thành một yêu cầu mua \\
Chiến lược & Học từ phần thưởng (RL) đối đầu luật cố định & Phân vai; luật xác minh ngân sách, tồn kho \\
Cấu trúc & Không giao tiếp; tương tác qua môi trường & Tập trung qua Coordinator \\
Điều phối & Đồng thời: cả hai chọn rồi môi trường áp dụng & Tĩnh: tìm $\to$ lọc, xếp hạng $\to$ kiểm kho \\
Lỗi lan truyền & Phản ứng giảm giá dây chuyền của đối thủ & Tồn kho cũ dẫn tới xác nhận sai \\\bottomrule
\end{tabular}}
\end{table}

\nhanxet{Hai hệ nằm ở hai đầu của khung: hệ thứ nhất cạnh tranh, không có thông điệp, mọi tương tác ẩn trong động
lực thị trường; hệ thứ hai hợp tác, mọi tương tác là thông điệp tường minh. Vì vậy chỉ số đánh giá cũng khác hẳn:
lợi nhuận cho hệ thứ nhất, độ đúng và chi phí phối hợp cho hệ thứ hai.}
